\documentclass[11pt]{article}
\usepackage[a4paper,margin=1in]{geometry}
\usepackage{amsmath,amssymb,amsthm}
\usepackage{graphicx}
\usepackage{xcolor}
\usepackage[hidelinks]{hyperref}
\usepackage{microtype}
\setlength{\emergencystretch}{2em}

\newtheorem{theorem}{Theorem}
\newtheorem{lemma}[theorem]{Lemma}
\newtheorem{example}[theorem]{Example}
\newcommand{\WA}{\mathrm{WA}}
\newcommand{\rankR}{\operatorname{rank}_{\mathbb R}}

\title{Weak amenability with one infinite-center rank-one factor}
\author{Automated mathematical discovery run \texttt{fa\_banach\_001}}
\date{13 August 2026}

\begin{document}
\maketitle

\begin{center}
\fcolorbox{black}{gray!8}{\parbox{0.91\linewidth}{
\textbf{Status.} Candidate substantial partial result, likely valid, subject
to human review.  The theorem proves the conjectured sufficiency direction
for connected Lie groups having at most one rank-one simple Levi factor with
infinite center.  It does not settle the remaining multi-factor central
mixing or the full necessity direction.
}}
\end{center}

\section{Source problem}

The source is S. Knudby, \emph{Weak amenability of Lie groups made
discrete}, arXiv:1612.04280~\cite{KnudbyDiscrete}.  On source p.~2 it recalls
that the weak-amenability classification for connected Lie groups was known
when the Levi factor has finite center and when the group is simply
connected, and then states: ``Some intermediate cases remain open.''  The
passage is reproduced below.

\begin{center}
\includegraphics[width=\linewidth]{figures/open_problem_crop.png}
\end{center}

More explicitly, write a Levi decomposition of the Lie algebra as
\[
 \mathfrak g=\mathfrak r\rtimes\mathfrak s,
 \qquad \mathfrak s=\bigoplus_{i=1}^k\mathfrak s_i,
\]
where $\mathfrak r$ is solvable and the $\mathfrak s_i$ are simple.  The
simply connected classification~\cite[Theorem~1.10]{KnudbySimplyConnected}
asserts that weak amenability is equivalent to
\begin{equation}\label{eq:C}
 \text{for every $i$, either }\rankR(\mathfrak s_i)=0,
 \quad\text{or}\quad
 \rankR(\mathfrak s_i)=1\ \text{ and }[\mathfrak s_i,\mathfrak r]=0.
\end{equation}
The open classification predicts the same criterion without a restriction
on global topology.

We write $\Lambda_{\WA}(H)$ for the Cowling--Haagerup constant.  We use the
standard permanence facts collected in~\cite[Section~2]{KnudbyDiscrete}:
compact normal quotients preserve this constant; direct products multiply
it; and if $N\triangleleft H$ is closed and $H/N$ is amenable, then
$\Lambda_{\WA}(N)=\Lambda_{\WA}(H)$.  We also use the source's semisimple
theorem~\cite[Theorem~3.1]{KnudbyDiscrete}, which in particular shows that
the constant of a connected simple Lie group is unchanged on passing among
its connected local forms.

\section{Proof intuition}

Under condition~\eqref{eq:C}, every noncompact Levi factor splits away from
the amenable part of the universal cover.  Factors with finite center can be
made adjoint by one finite compact quotient.  What remains is an amenable
group $A$ and, at worst, one simple group $U$ whose center is infinite
cyclic, modulo an arbitrary discrete central subgroup $D$.

The key is to close the image of $U$.  Its quotient in the ambient group is
amenable, so it alone determines the Cowling--Haagerup constant.  Modulo the
part of $D$ already lying in $A$, the kernel is cyclic.  A cyclic subgroup
of an abelian Lie group either is closed (unless its closure is compact) or
has compact closure.  In the closed case the closure of $U$ is just $U$;
in the compact case it differs from the adjoint form $U/Z(U)$ only by a
compact central group.  Both alternatives have exactly the expected
constant.

Two infinite cyclic simple centers are qualitatively different: a
rank-two kernel can have dense image in a noncompact real direction.  This
explains both the reach and the stopping point of the argument.

\section{The cyclic-closure lemma}

\begin{lemma}\label{lem:cyclic}
Let $A$ be an amenable Lie group and let $U$ be a connected simple Lie group
whose center is infinite cyclic.  Let
\[
 D\leq Z(A)\times Z(U)
\]
be discrete, and put $Q=(A\times U)/D$.  Then
\[
 \Lambda_{\WA}(Q)=\Lambda_{\WA}(U).
\]
\end{lemma}

\begin{proof}
Let $q:A\times U\to Q$ be the quotient map and set
$N=\overline{q(U)}$.  The group $N$ is closed and normal, and $Q/N$ is a
quotient of $A$, hence amenable.  Therefore
\begin{equation}\label{eq:coamenable}
 \Lambda_{\WA}(Q)=\Lambda_{\WA}(N).
\end{equation}

Fix a generator $z$ of $Z(U)$, let
\[
 A_0=\overline{\operatorname{pr}_A(D)},
 \qquad D_A=D\cap(A\times\{e\}).
\]
Because multiplying an element of $D$ by $U$ absorbs its $U$-coordinate,
\[
 q^{-1}(N)=\overline{UD}=A_0\times U,
 \qquad N=(A_0\times U)/D.
\]
The group $D_A$ is discrete central in $A_0$.  If the projection of $D$ to
$Z(U)$ is trivial, then $D=D_A$, $A_0=D_A$, and $N\cong U$, proving the
claim.

Otherwise this projection is $\langle z^\ell\rangle$ for some $\ell\geq1$.
Choose $(a,z^\ell)\in D$.  Then $D/D_A$ is generated by the class of this
element.  In the abelian Lie group
\[
 B=A_0/D_A,
\]
the cyclic subgroup generated by $b=aD_A$ is dense by the definition of
$A_0$.  The elementary structure theorem for abelian Lie groups now gives
the following dichotomy: either $B$ is compact, or $B=\langle b\rangle$ is
an infinite discrete cyclic group.  Indeed, a dense cyclic subgroup cannot
be dense in a nonzero Euclidean factor; if its projection to a discrete
factor is nonzero, convergence forces the exponent to be eventually
constant, making the cyclic subgroup closed; with neither type of factor,
the closure is compact.

If $B=\langle b\rangle$ is infinite discrete, every element of $A_0$ can be
absorbed using a power of $(a,z^\ell)$ and an element of $D_A$.  Thus the
map $U\to N$ is surjective.  Its kernel is trivial, since $b$ has infinite
order, so $N\cong U$.

Suppose instead that $B$ is compact.  After dividing by $D_A$ we may write
\[
 N\cong (B\times U)/\langle(b,z^\ell)\rangle.
\]
The image
\[
 C=(B\times Z(U))/\langle(b,z^\ell)\rangle
\]
is central in $N$ and is compact: every one of its elements has a
representative in
\[
 B\times\{e,z,\ldots,z^{\ell-1}\},
\]
a compact set.  Moreover $N/C\cong U/Z(U)$.  Compact-quotient invariance and
the semisimple theorem therefore yield
\[
 \Lambda_{\WA}(N)=\Lambda_{\WA}(U/Z(U))
 =\Lambda_{\WA}(U).
\]
Together with~\eqref{eq:coamenable}, this proves the lemma.
\end{proof}

\section{Partial classification theorem}

Among the real-rank-one simple real Lie algebras, precisely those of type
$\mathfrak{su}(n,1)$, $n\geq1$, have a simply connected Lie group with
infinite center; here $\mathfrak{su}(1,1)\cong\mathfrak{sl}(2,\mathbb R)$.

\begin{theorem}\label{thm:main}
Let $G$ be a connected Lie group whose Lie algebra satisfies
condition~\eqref{eq:C}.  Suppose that at most one rank-one simple ideal
$\mathfrak s_i$ is of type $\mathfrak{su}(n,1)$ for some $n\geq1$.
Then $G$ is weakly amenable.  More precisely, if $\widetilde S_i$ is the
simply connected simple Lie group with Lie algebra $\mathfrak s_i$, then
\[
 \Lambda_{\WA}(G)=\prod_{i=1}^k
 \Lambda_{\WA}(\widetilde S_i).
\]
Thus the conjectured sufficiency direction and its predicted constant hold
for this class.
\end{theorem}

\begin{proof}
Let $\widetilde G\to G$ be the universal covering homomorphism and let $D$
be its discrete central kernel.  In the simply connected Levi decomposition,
separate the compact simple factors from the noncompact ones.  By
condition~\eqref{eq:C}, every noncompact factor commutes with the solvable
radical and therefore is a direct factor.  Consequently
\[
 \widetilde G\cong A\times P\times U.
\]
Here $A$ is the semidirect product of the simply connected solvable radical
by the compact Levi factors, and is amenable; $P$ is the product of the
rank-one factors having finite center; and $U$ is absent or is the unique
rank-one factor having infinite cyclic center.

Let $F=Z(P)$, a finite group, and let $K=(DF)/D\triangleleft G$.  The group
$K$ is finite, hence compact.  Passing to $G/K$ does not change the
Cowling--Haagerup constant.  In $\widetilde G/F$, the factor $P/F$ is
center-free.  Therefore the image $D'$ of $DF/F$ has trivial $P/F$
coordinate and
\[
 G/K\cong (P/F)\times Q,
 \qquad Q=(A\times U)/D',
 \qquad D'\leq Z(A)\times Z(U).
\]
The group $D'$ is still discrete, since only a finite group was removed.

If $U$ is absent, then $Q$ is amenable.  If $U$ is present,
Lemma~\ref{lem:cyclic} gives
$\Lambda_{\WA}(Q)=\Lambda_{\WA}(U)$.  Hence product and compact-quotient
invariance give
\[
 \Lambda_{\WA}(G)
 =\Lambda_{\WA}(P/F)\Lambda_{\WA}(Q)
 =\prod_{i=1}^k\Lambda_{\WA}(\widetilde S_i).
\]
The last equality uses the semisimple local-form theorem and the fact that
compact factors have constant one.  In particular the displayed value is
finite, so $G$ is weakly amenable.
\end{proof}

\begin{example}[A genuinely intermediate topology]\label{ex:irrational}
Let $U$ be the universal cover of $\mathrm{SU}(n,1)$, let $z$ generate
$Z(U)$, and fix $\alpha\notin\mathbb Q$.  In $\mathbb R\times U$ set
\[
 D=\langle(1,e),(\alpha,z)\rangle,
 \qquad G=(\mathbb R\times U)/D.
\]
The subgroup $D$ is discrete although its projection
$\mathbb Z+\alpha\mathbb Z$ to $\mathbb R$ is dense.  The group $G$ is
neither simply connected nor covered by the finite-center-Levi hypothesis;
indeed the analytic image of $U$ is injective and dense.  Nevertheless
Theorem~\ref{thm:main} gives
$\Lambda_{\WA}(G)=\Lambda_{\WA}(U)$.  This shows that the theorem is not
merely a reformulation of either previously known extreme case.
\end{example}

\section{Upgrade attempts and remaining obstruction}

Eight focused routes were audited.  First, lifting an arbitrary lattice to
the universal cover does \emph{not} descend weak amenability: Cowling's
discrete central-quotient inequality has the opposite direction from the
one needed.  Second, a lattice that splits in the universal cover would
work, as in the source's free-lattice argument for $\mathrm{SL}(2,\mathbb
R)$, but splitting is a real cohomological obstruction for complex
hyperbolic lattices.  Third, periodizing compactly supported Herz--Schur
multipliers over $D$ has no known uniform completely bounded norm estimate.

Fourth, the cyclic-closure route gives Lemma~\ref{lem:cyclic} and
Theorem~\ref{thm:main}.  Fifth, it cannot simply be iterated: with two
infinite centers the homomorphism
\[
 \mathbb Z^2\longrightarrow\mathbb R,
 \qquad(m,n)\longmapsto m+\alpha n
\]
can have dense noncompact image, neither compact nor discrete.  Sixth, an
attempt at necessity via locally bad analytic subgroups stops because such
subgroups need not be closed in intermediate topology.  Seventh, no
non-weakly-amenable closed discrete subgroup was obtained from central
quotients of rank-one lattices.  Eighth, the later-literature search
described below found no theorem removing the remaining obstruction.

Thus the packet makes no claim when two or more
$\mathfrak{su}(n,1)$ factors can mix simultaneously with a noncompact
central direction.  It also does not claim the necessity half in arbitrary
intermediate topology.  The source proves additional special cases by
different lattice methods, notably multiple $\mathfrak{sl}_2$ factors at
constant one; the present result is complementary, not a maximal catalogue.

\section{Novelty audit and verification recommendation}

A bounded search on 13 August 2026 covered arXiv:1612.04280 and
arXiv:1505.00984, the exact phrase ``some intermediate cases remain open'',
``weak amenability connected Lie groups arbitrary fundamental group'',
``reductive Lie groups lattices'', and ``lattice lifts to the universal
cover'' for $\mathrm{SU}(n,1)$.  It found the known finite-center,
simply-connected, semisimple, and special $\mathrm{SL}(2,\mathbb R)$
results, plus later work on residual finiteness of complex-hyperbolic
central extensions, but no full later classification or this cyclic-closure
subcase.  Novelty confidence is provisional.

Human review should focus on two points: the decomposition after killing
$Z(P)$ in Theorem~\ref{thm:main}, and the claim that the dense cyclic group
in Lemma~\ref{lem:cyclic} has either compact closure or is itself closed.
Both have been expanded in the proof so that no multiplier construction or
lattice-splitting assumption is hidden.

\begin{thebibliography}{9}
\bibitem{KnudbyDiscrete}
S. Knudby,
\emph{Weak amenability of Lie groups made discrete},
Pacific J. Math. 297 (2018), 101--116;
arXiv:1612.04280.

\bibitem{KnudbySimplyConnected}
S. Knudby,
\emph{Weak amenability and simply connected Lie groups},
Kyoto J. Math. 56 (2016), 689--700;
arXiv:1505.00984.

\bibitem{CDSW}
M. Cowling, B. Dorofaeff, A. Seeger, and J. Wright,
\emph{A family of singular oscillatory integral operators and failure of
weak amenability},
Duke Math. J. 127 (2005), 429--486.

\bibitem{CowlingHaagerup}
M. Cowling and U. Haagerup,
\emph{Completely bounded multipliers of the Fourier algebra of a simple Lie
group of real rank one},
Invent. Math. 96 (1989), 507--549.
\end{thebibliography}

\end{document}
